
\begin{frame}
	\frametitle{Matematika, 1. del}
 Analiza, logika in množice.
	
  \end{frame}

\begin{frame}{Logika in množice}
	\begin{enumerate}
		\item
		Poišči preneksno obliko formule \[
\exists x : P(x) \land \forall x : Q(x) \Rightarrow \forall x : R(x)
\].
		\item 
		Definiramo množici $A$=[2,5] in $B$={0,1,2,3,4\dots}.
		V ravnino nariši:
		\begin{enumerate}
		   \item \[ \displaystyle
A \cap B \times \emptyset
\]
		   \item \[ \displaystyle
(A \cup B)^c \times \mathbb{R}
\]
		\end{enumerate}
		\item
		Dokaži:
		\begin{itemize}
			\item Prva enačba:
			\begin{align*}
				\neg (A \lor B) &\sim (\neg A \land \neg B)
			\end{align*}
			\item Druga enačba:
			\begin{align*}
				(A \Rightarrow B) &\sim (\neg B \Rightarrow \neg A)
			\end{align*}
			\item Tretja enačba:
			\begin{align*}
				A \cap B &\times \emptyset
			\end{align*}
	    \end{itemize}
	\end{enumerate}
\end{frame}

\begin{frame}{Analiza}
	\begin{enumerate}
		\item
		Pokaži, da je funkcija \[ x\to\sqrt{x} \] enakomerno zvezna na ??.
		\item 
		Katero krivuljo določa sledeč parametričen zapis?
		% Spodaj si pomagajte z dokumentacijo o razmikih v matematičnem načinu.
		% https://www.overleaf.com/learn/latex/Spacing_in_math_mode
		$$
		   x(t) = a \cos t, ?? % tu manjka ukaz za presledek
		   y(t) = b \sin t, ?? % tu manjka ukaz za presledek
		   t \in [0, 2 \pi]
		$$ 
		\item
		Pokaži, da ima ?? inverzno funkcijo in izračunaj ??.
		
		\item
		Izračunaj integral 
		% V rešitvah smo spodnji integral zapisali v vrstičnem načinu,
		% ampak v prikaznem slogu. To naredite tako, da v matematičnem načinu najprej
		% uporabite ukaz displaystyle.
		% Pred dx je presledek: pravi ukaz je \,
		??
		% \frac{2+\sqrt{x+1}}{(x+1)^2-\sqrt{x+1}} 
		\item 
		Naj bo $g$ zvezna funkcija. Ali posplošeni integral 
		??
		konvergira ali divergira? Utemelji.
	\end{enumerate}
\end{frame}

\begin{frame}{Kompleksna števila}
	\begin{enumerate}
		\item
		Naj bo $z$ kompleksno število, $z \ne 1$ in ??.
		Dokaži, da je število \( i \, \frac{z+1}{z-1} \) realno.
		\item
		Poenostavi izraz:
		??
	\end{enumerate}
\end{frame}